% Part: turing-machines
% Chapter: undecidability
% Section: trakhtenbrot

\documentclass[../../../include/open-logic-section]{subfiles}

\begin{document}

\olfileid{tur}{und}{tra}
\olsection{ટ્રાખ્ટેનબ્રોટનું પ્રમેય}

\begin{explain}
  \olref[rep]{sec}માં આપણે ટ્યુરિંગ મશીન~$M$
  અને ઇનપુટ શબ્દમાળા~$w$ માટે !!{sentence}s
  $!T(M,w)$ તથા~$!E(M,w)$ વ્યાખ્યાયિત કર્યા.
  પછી \olref[ver]{lem:valid-if-halt} અને
  \olref[ver]{lem:halt-if-valid}માં બતાવ્યું કે
  $!T(M,w) \lif !E(M,w)$ ત્યારે અને માત્ર ત્યારે
  પ્રામાણ્ય ધરાવે છે જ્યારે ઇનપુટ~$w$ પર શરૂ થયેલું
  $M$ છેવટે અટકે છે. અટકવાની સમસ્યા અનિર્ણેય
  હોવાથી, આમાંથી પ્રથમ-ક્રમના તર્કનાં !!{sentence}sનું
  પ્રામાણ્ય અને સંતોષ્યતા અનિર્ણેય છે એ ફલિત થાય છે
  (\cref{tur:und:uns:thm:decision-prob,tur:und:uns:cor:undecidable-sat}).

  પરંતુ વાક્યનું પ્રામાણ્ય અને સંતોષ્યતા કોઈ પણ
  !!{structure}s માટે વ્યાખ્યાયિત છે—સાન્ત હોય કે
  અનંત. તમને એવું લાગે કે કોઈ !!a{sentence}
  સાન્ત !!{structure}માં સંતોષ્ય છે કે નહીં (અથવા
  બધાં સાન્ત !!{structure}sમાં પ્રામાણ્ય ધરાવે છે કે
  નહીં) તે નક્કી કરવું સરળ હશે. નિર્ણય સમસ્યા
  ઉકેલી ન શકાય તેની સાબિતી અનુરૂપ ફેરવીને
  બતાવી શકીએ કે એવું નથી.

  પહેલાં \olref[ver]{lem:halt-if-valid}ની
  સાબિતી ફરી જુઓ: ત્યાં આપણે~$!T(M,w)$નું
  નિદર્શ~$\Struct M$ બનાવ્યું હતું, જે
  ઇનપુટ~$w$ પર શરૂ થયેલું મશીન~$M$ શું કરે
  છે તેનું ચોક્કસ વર્ણન કરે છે. એ નિદર્શનો
  પ્રદેશ~$\Nat$ હોવાથી તે અનંત હતો. પરંતુ
  $M$ ઇનપુટ~$w$ પર ખરેખર અટકતું હોય, તો
  એ જ રીતે સાન્ત નિદર્શ~$\Struct M'$ બનાવી શકીએ.
  માનો કે ઇનપુટ~$w$ પર શરૂ થયેલું~$M$,
  $k$ પગલાં પછી અટકે છે. પ્રદેશ
  $\Domain{M'}$ તરીકે $\{0, \dots, n\}$ લો,
  જ્યાં~$n$ એ~$k+1$ અને~$w$ની લંબાઈમાંથી
  મોટું છે, અને મૂકો:
  \sourcecorrection{OLUN-013}{સ્થિર અંગ્રેજી મૂળ
  $n$ માટે $k$ અને ઇનપુટની લંબાઈમાંથી મોટું
  મૂલ્ય લે છે. જો $n=k$ હોય, તો અટકવાના
  સમયનો પૂર્વગામી સંક્રમણ અથવા છેલ્લે લખાતું
  ખાનું સ્વ-લૂપના છેડે આવી શકે છે. $n>k$
  રાખવાથી વપરાતી સમય-સંખ્યા અને અટકવા પહેલાં
  લખાતાં ખાનાં એ છેડાથી નીચે રહે છે.}
  \[
    \Assign{\prime}{M'}(x) =
    \begin{cases}
      x + 1 &\text{જો $x < n$ હોય}\\
      n &\text{અન્યથા,}
    \end{cases}
  \]
  અને $\tuple{x,y} \in \Assign{<}{M'}$ ત્યારે
  અને માત્ર ત્યારે જ્યારે $x < y$ અથવા
  $x = y = n$. બાકીની રીતે $\Struct{M'}$ની
  વ્યાખ્યા~$\Struct{M}$ જેવી જ છે.
  $\Struct{M'}$ની વ્યાખ્યા મુજબ,
  \olref[ver]{lem:halt-if-valid}ની સાબિતીની
  જેમ, $\Sat{M'}{!T(M,w)}$. અને આપણે~$M$
  ઇનપુટ~$w$ પર અટકે છે એમ માન્યું હોવાથી
  $\Sat{M'}{!E(M,w)}$. તેથી
  $\Struct{M'}$ એ~$!T(M,w) \land !E(M,w)$નું
  સાન્ત નિદર્શ છે (અહીં~$\lif$ને બદલે~$\land$
  મૂક્યું છે).

  આ સાબિતીનો અડધો ભાગ થયો: $M$ ઇનપુટ~$w$
  પર અટકે, તો $!T(M,w) \land !E(M,w)$ને
  સાન્ત નિદર્શ છે. દુર્ભાગ્યે ઊલટી દિશા સાચી
  નથી: કેટલાંક ટ્યુરિંગ મશીનો ઇનપુટ~$w$ પર
  અટકતાં નથી, છતાં $!T(M,w) \land !E(M,w)$ને
  સાન્ત નિદર્શ મળે છે. દાખલા તરીકે, એક જ
  અવસ્થા~$q_0$ અને સૂચના
  $\delta(q_0,\TMblank) = \tuple{q_0,\TMblank,\TMstay}$
  ધરાવતું મશીન~$M$ લો. ખાલી ઇનપુટ
  $w = \emptyseq$ પર તે કદી અટકતું નથી:
  તે અનંત ચક્રમાં ચાલે છે, પરંતુ ટેપ બદલતું કે
  હેડ ખસેડતું નથી. બધાં સંરૂપણો સરખાં જ છે
  (એક જ અવસ્થા, હેડનું સ્થાન અને ટેપની સામગ્રી).
  આપણે $!T(M,\emptyseq) \land !E(M,\emptyseq)$
  સંતોષે એવું સાન્ત !!{structure}~$\Struct{M''}$
  વ્યાખ્યાયિત કરી શકીએ (કસરત). તેમ છતાં,
  આવાં !!{structure}s દૂર થાય તે રીતે
  $!T(M,w)$માં યોગ્ય ફેરફાર કરી શકીએ.

  \begin{prob}
    એક જ અવસ્થા~$q_0$ અને એક જ સૂચના
    $\delta(q_0,\TMblank) = \tuple{q_0,\TMblank,\TMstay}$
    ધરાવતું ટ્યુરિંગ મશીન~$M$ લો.
    $\Domain{M''} = \{0, 1, 2\}$,
    $\Assign{\prime}{M''}(0) =
    \Assign{\prime}{M''}(1) = 1$ અને
    $\Assign{\prime}{M''}(2) = 2$, તથા
    $\Assign{<}{M''} = \{\tuple{0,1}, \tuple{1,1}, \tuple{2,2}\}$
    લો.
    \sourcecorrection{OLUN-015}{સ્થિર અંગ્રેજી
    મૂળ ``એક જ અવસ્થા'' કહેતાં છતાં સૂચનામાં
    અપરિચિત $q$ આપે છે; અનુગામીની બીજી
    કિંમતમાં $M''$ને બદલે $M'$ છે. અહીં
    બંને જગ્યાએ આ જ મશીન અને નિદર્શનાં
    યોગ્ય ચિહ્નો રાખ્યાં છે.}
    $\Assign{\Obj Q_{q_0}}{M''}$,
    $\Assign{\Obj S_{\TMblank}}{M''}$ અને
    $\Assign{\Obj S_{\TMendtape}}{M''}$
    એ રીતે વ્યાખ્યાયિત કરો કે
    $!T(M,\emptyseq)$ તથા $!E(M,\emptyseq)$
    સાચાં થાય, અને કેમ થાય છે તે સમજાવો.
    સંકેત: $\delta(q_0, \TMendtape)$
    અવ્યાખ્યાયિત છે એ નોંધો. ખાતરી કરો કે
    \begin{align*}
    & \Obj Q_{q_0}(\num{1}, \num{n}) \land \Obj S_{\TMendtape}(\num{0}, \num{n})
      \land
      \lforall[x][(\num{0} < x \lif \Obj S_\TMblank(x, \num{n}))] \text{\quad બધાં $n \in \Nat$ માટે}\\
    & \lexists[y][(\Obj Q_{q_0}(\num{0}, y) \land \Obj S_{\TMendtape}(\num{0}, y))]
    \end{align*}
    બંને~$\Struct{M''}$માં સત્ય છે.
  \end{prob}
\end{explain}

ઇનપુટ $w = \sigma_{i_1}\dots\sigma_{i_k}$ પર
ટ્યુરિંગ મશીન~$M$નું કામ વર્ણવતાં !!{sentence}s
વિચારો:
\begin{enumerate}
\item સંખ્યાઓ અને~$<$ વર્ણવતાં સ્વયંસિદ્ધ વાક્યો
  (\olref[rep]{sec}માં $!T(M,w)$ની વ્યાખ્યા જેવા જ).
\item પ્રારંભિક સંરૂપણ વર્ણવતાં સ્વયંસિદ્ધ વાક્યો:
  $!T(M,w)$ની વ્યાખ્યા જેવા જ.
\item એક સંરૂપણમાંથી બીજા સંરૂપણમાં થતું સંક્રમણ
  વર્ણવતાં સ્વયંસિદ્ધ વાક્યો:

આગળ માટે $!A(x, y)$ અગાઉ જેવું જ માનો અને
નીચેનું વ્યાખ્યાયિત કરો:
\[
  !B(y) \ident \lforall[x][(x < y \lif \eq/[x][y])].
\]
\begin{enumerate}
\item \ollabel{rep-right} દરેક સૂચના
  $\delta(q_i, \sigma) =
  \tuple{q_j, \sigma', \TMright}$ માટે !!{sentence}:
\begin{align*}
& \lforall[x][\lforall[y][(
   (\Obj Q_{q_i}(x, y) \land \Obj S_{\sigma}(x, y)) \lif {}]] \\
&\qquad   (\Obj Q_{q_j}(x', y') \land \Obj S_{\sigma'}(x, y') \land
!A(x, y) \land !B(y')))
\end{align*}
\item \ollabel{rep-left} દરેક સૂચના
  $\delta(q_i, \sigma) =
  \tuple{q_j, \sigma', \TMleft}$ માટે !!{sentence}:
\begin{align*}
& \lforall[x][\lforall[y][
    ((\Obj Q_{q_i}(x', y) \land \Obj S_{\sigma}(x', y)) \lif {}]]\\
& \qquad   (\Obj Q_{q_j}(x, y') \land \Obj S_{\sigma'}(x', y') \land
!A(x', y) \land !B(y'))) \land {}\\
& \lforall[y][((\Obj Q_{q_i}(\num{0}, y) \land \Obj S_{\sigma}(\num{0},
    y)) \lif {}]\\
& \qquad (\Obj Q_{q_j}(\num{0}, y') \land \Obj S_{\sigma'}(\num{0},
  y') \land !A(\num{0}, y) \land !B(y')))
\end{align*}
\sourcecorrection{OLUN-016}{સ્થિર અંગ્રેજી મૂળની
સામાન્ય ડાબે-ચાલ શાખામાં લખાતું ખાનું~$x'$
છતાં $!A(x,y)$ છે અને આગળના સમય માટે
$!B(y')$ નથી. અહીં લખાતું ખાનું
યથાવત્-શરતમાંથી બાકાત છે અને આ શાખામાં પણ
નવો સમય અગાઉના સમયોથી જુદો રહે છે.}
\item \ollabel{rep-stay} દરેક સૂચના
  $\delta(q_i, \sigma) =
  \tuple{q_j, \sigma', \TMstay}$ માટે !!{sentence}:
\begin{align*}
& \lforall[x][\lforall[y][(
   (\Obj Q_{q_i}(x, y) \land \Obj S_{\sigma}(x, y)) \lif {}]] \\
&\qquad (\Obj Q_{q_j}(x, y') \land \Obj S_{\sigma'}(x, y') \land
!A(x, y) \land !B(y')))
\end{align*}
\end{enumerate}
જેમ જોઈ શકાય છે,~$M$નાં સંક્રમણો વર્ણવતાં
!!{sentence}s, અંતે $!B(y')$ ઉમેર્યું છે તે સિવાય,
$!T(M,w)$નાં અનુરૂપ !!{sentence} જેવા જ છે.
$!B(y')$ ખાતરી કરે છે કે ``આગળના'' સંરૂપણનો
સૂચક~$y'$ અગાઉની બધી સંખ્યાઓ
$\Obj 0$, $\Obj 0'$, \dots{}થી જુદો છે.
% \item સુસંગતતાની શરતો વ્યક્ત કરતાં વાક્યો:
% \begin{enumerate}
%   \item\ollabel{rep-con-symb}%
%   કોઈ પણ સમયે એક ખાનામાં એક કરતાં વધુ ચિહ્ન ન હોઈ શકે એવું
%   !!^a{sentence}:
%   \[\lforall[x][\lforall[y][(\Obj S_{\sigma}(x, y) \lif \lnot \Obj
%   S_{\sigma'}(x, y))]],\]
%   $T$ના ટેપ-ચિહ્નોની દરેક જોડી $\sigma \neq \sigma'$ માટે.
%   \item\ollabel{rep-con-state}%
%   કોઈ પણ સમયે $M$ એક કરતાં વધુ અવસ્થામાં ન હોઈ શકે એવું
%   !!^a{sentence}:
%   \[\lforall[x][\lforall[y][(\Obj Q_{q}(x, y) \lif
%   \lforall[z][\lnot \Obj
%   Q_{q'}(z, y))]],\]
%   $T$ની અવસ્થાઓની દરેક જોડી $q \neq q'$ માટે.
%   \item\ollabel{rep-con-tape}%
%   કોઈ પણ સમયે $M$ એક કરતાં વધુ ટેપ-ખાનાં પર ન હોઈ શકે એવું
%   !!^a{sentence}:
%   \[\lforall[x][\lforall[y][(\Obj Q_{q}(x, y) \lif
%   \lforall[z][\eq/[z][y]\lif \lnot \Obj
%   Q_{q'}(z, y))]],\]
%   $T$ની અવસ્થાઓની દરેક જોડી $q$, $q'$ માટે.
% \end{enumerate}
\end{enumerate}
ટ્યુરિંગ મશીન~$M$ અને ઇનપુટ~$w$ માટે ઉપરનાં
બધાં !!{sentence}sનું સંયોજન~$!T'(M, w)$ માનો.

\begin{lem}\ollabel{lem:halts-sat}
  ઇનપુટ~$w$ પર શરૂ થયેલું~$M$ અટકે, તો
  $!T'(M,w) \land !E(M,w)$ને સાન્ત નિદર્શ છે.
\end{lem}

\begin{proof}
  $\Struct{M'}$ને
  \olref[ver]{lem:halt-if-valid}ની સાબિતી જેવું
  જ લો, પરંતુ:
  \begin{align*}
    \Domain{M'} & = \{0, \dots, n\},\\
    \Assign{\prime}{M'}(x) & =
    \begin{cases}
      x + 1 &\text{જો $x < n$ હોય}\\
      n &\text{અન્યથા,}
    \end{cases} \\
    \tuple{x,y} \in \Assign{<}{M'} &\text{ત્યારે અને માત્ર ત્યારે જ્યારે $x < y$ અથવા $x = y = n$,}
  \end{align*}
  જ્યાં $n = \max(k+1,\len{w})$ અને~$k$ એ
  એવી સૌથી નાની પગલાંની સંખ્યા છે કે
  ઇનપુટ~$w$ પર શરૂ થયેલું~$M$, $k$ પગલાં
  પછી અટક્યું હોય.
  \sourcecorrection{OLUN-014}{સ્થિર અંગ્રેજી
  મૂળ અહીં ફરી $n=\max(k,\len{w})$ લે છે.
  OLUN-013 મુજબ અટકવાના સમય અને છેલ્લા
  લખાતા ખાનાને છેડાની સ્વ-લૂપથી દૂર રાખવા
  અહીં $k+1$ લીધું છે; એટલે $!B(\num{k})$
  પણ અંતિમ સંક્રમણથી ફલિત થઈ શકે છે.}
  $\Sat{M'}{!T'(M,w) \land !E(M,w)}$
  તેની ચકાસણી કસરત તરીકે છોડી છે.
  \sourcecorrection{OLUN-017}{સ્થિર અંગ્રેજી
  મૂળના આ સંતોષ-દાવામાં $E(M,w)$ આગળનું
  મેટાસંકેત~$!$ છૂટ્યું છે; અહીં વ્યાખ્યાયિત
  વાક્ય~$!E(M,w)$ જ રાખ્યું છે.}
\end{proof}

\begin{prob}
  $\Sat{M'}{!T'(M,w) \land !E(M,w)}$
  સાબિત કરીને \olref[tur][und][tra]{lem:halts-sat}ની
  સાબિતી પૂર્ણ કરો.
  \sourcecorrection{OLUN-018}{સ્થિર અંગ્રેજી
  મૂળની આ કસરતમાં $!T'(M,w)$ને બદલે
  $!T(M,w)$ છે અને $E(M,w)$ આગળનું $!$
  છૂટે છે; ઉપપ્રમેયમાં દર્શાવેલું જ સંયોજન
  અહીં ચકાસવાનું છે.}
\end{prob}

\begin{lem}\ollabel{lem:sat-halts}
  $!T'(M,w) \land !E(M,w)$ને સાન્ત નિદર્શ
  હોય, તો ઇનપુટ~$w$ પર શરૂ થયેલું~$M$ અટકે છે.
\end{lem}

\begin{proof}
  પ્રતિનિષેધ સાબિત કરીએ. માનો કે ઇનપુટ~$w$
  પર શરૂ થયેલું~$M$ અટકતું નથી.
  $!T'(M,w) \land !E(M,w)$ને કોઈ નિદર્શ
  જ ન હોય, તો વાત પૂરી. તેથી માનો કે
  $\Struct{M}$ એ~$!T'(M,w) \land !E(M, w)$નું
  નિદર્શ છે. હવે બતાવવું છે કે તે સાન્ત હોઈ
  શકતું નથી.
  \sourcecorrection{OLUN-019}{સ્થિર અંગ્રેજી
  મૂળ અહીં જૂના~$!T(M,w)$નું નિદર્શ માને છે,
  પણ નવી~$!B$ શરત પરથી અનંતપણું બતાવવા
  $!T'(M,w)$નું નિદર્શ જરૂરી છે.}

  \olref[ver]{lem:config}ની જેમ સાબિત કરી
  શકીએ કે ઇનપુટ~$w$ પર શરૂ થયેલું~$M$ જો
  $n$ પગલાં પહેલાં અટક્યું ન હોય અને
  $n \ge 1$ હોય, તો $!T'(M,w) \Entails
  !C(M, w, n) \land !B(\num{n})$.
  $M$ ઇનપુટ~$w$ પર અટકતું નથી, તેથી
  દરેક ધન~$n \in \Nat$ માટે
  $!T'(M,w) \Entails !C(M, w, n) \land
  !B(\num{n})$. પ્રારંભિક સંરૂપણનું
  $!C(M,w,0)$ ઇનપુટનાં સ્વયંસિદ્ધ વાક્યોમાંથી
  અલગથી મળે છે.
  \sourcecorrection{OLUN-020}{સ્થિર અંગ્રેજી
  મૂળ $!B(\num{n})$ને $n=0$ માટે પણ
  ફલિત કહે છે; પ્રારંભિક સ્વયંસિદ્ધ વાક્યો
  એ દાવો આપતાં નથી. પ્રથમ સંક્રમણથી
  આગળના બધા $n\ge1$ માટે~$!B$ મળે છે,
  અને જુદા આંકિક પદો બતાવવા એટલું પૂરતું છે.}
  \olref[rep]{prop:mlessk}થી, દરેક~$k < n$
  માટે $!T'(M,w) \Entails \num{k} < \num{n}$.
  વળી $!B(\num{n}) \Entails
  \num{k} < \num{n} \lif
  \eq/[\num{k}][\num{n}]$. તેથી બધા
  $k < n$ માટે
  $\Sat{M}{\eq/[\num{k}][\num{n}]}$:
  એટલે અસંખ્ય પદો~$\num{k}$નાં મૂલ્યો
  $\Struct{M}$માં જુદાં જ હોવાં જોઈએ.
  પરંતુ તેના માટે $\Domain{M}$ અનંત હોવો
  જરૂરી છે. તેથી~$\Struct{M}$,
  $!T'(M,w) \land !E(M, w)$નું સાન્ત નિદર્શ
  હોઈ શકતું નથી.
\end{proof}

\begin{prob}
  ઇનપુટ~$w$ પર શરૂ થયેલું~$M$ જો~$n$
  પગલાં પહેલાં અટક્યું ન હોય અને~$n \ge 1$,
  તો $!T'(M,w) \Entails !B(\num{n})$
  સાબિત કરીને
  \olref[tur][und][tra]{lem:sat-halts}ની સાબિતી
  પૂર્ણ કરો.
  \sourcecorrection{OLUN-021}{સ્થિર અંગ્રેજી
  મૂળ આ કસરતમાં $n=0$ પણ આવરી લે છે.
  OLUN-020 મુજબ તે પ્રારંભિક કિસ્સામાં
  આ દાવો જરૂરી નથી કે ફલિત પણ થતો નથી.}
\end{prob}

\begin{thm}[ટ્રાખ્ટેનબ્રોટનું પ્રમેય]
  \ollabel{thm:trakhtenbrodt}
  પ્રથમ-ક્રમના તર્કનું મનસ્વી !!{sentence} સાન્ત
  નિદર્શ ધરાવે છે કે નહીં (એટલે સાન્ત
  નિદર્શમાં સંતોષ્ય છે કે નહીં) તે અનિર્ણેય છે.
\end{thm}

\begin{proof}
  માનો કે સાન્ત નિદર્શમાં સંતોષ્યતા નક્કી કરતું
  ટ્યુરિંગ મશીન~$F$ હોય. તો કોઈ પણ ટ્યુરિંગ
  મશીન~$M$ અને ઇનપુટ~$w$ માટે વાક્ય
  $!T'(M,w) \land !E(M,w)$ ગણીને~$F$ વડે
  તેને સાન્ત નિદર્શ છે કે નહીં તે નક્કી કરી
  શકીએ. \cref{tur:und:tra:lem:halts-sat,tur:und:tra:lem:sat-halts}
  મુજબ એવું નિદર્શ ત્યારે અને માત્ર ત્યારે છે
  જ્યારે ઇનપુટ~$w$ પર શરૂ થયેલું~$M$ અટકે.
  એટલે~$F$ વડે અટકવાની સમસ્યા ઉકેલી શકાય;
  પરંતુ તે ઉકેલી શકાતી નથી એ જાણીએ છીએ.
\end{proof}

\begin{cor}\ollabel{cor:fproof-incomp}%
  \emph{સાન્ત} નિદર્શોમાં પ્રામાણ્ય માટે
  યથાર્થ અને પૂર્ણ એવી કોઈ !!{derivation}
  પદ્ધતિ હોઈ શકતી નથી: એટલે એવી !!a{derivation} પદ્ધતિ,
  જેમાં દરેક સાન્ત !!{structure}~$\Struct{M}$
  માટે $\Sat{M}{!B}$ હોય ત્યારે અને માત્ર
  ત્યારે $\Proves !B$ હોય.
\end{cor}

\begin{proof}
  કસરત.
\end{proof}

\begin{prob}
  \olref[tur][und][tra]{cor:fproof-incomp}
  સાબિત કરો. $!B$ દરેક સાન્ત !!{structure}માં
  સંતોષાય ત્યારે અને માત્ર ત્યારે~$\lnot !B$
  સાન્ત નિદર્શમાં સંતોષ્ય નથી એ નોંધો.
  \olref[tur][und][uns]{thm:valid-ce}ના અર્થમાં
  સાન્ત નિદર્શમાં સંતોષ્યતા અર્ધનિર્ણેય કેમ
  છે તે સમજાવો. આથી સાન્ત નિદર્શોમાં
  પ્રામાણ્ય માટેની !!a{derivation} પદ્ધતિ હોત,
  તો સાન્ત નિદર્શમાં સંતોષ્યતા નિર્ણેય બનત
  એવું કેમ થાય તે દલીલ કરો.
\end{prob}

\end{document}
